\section{A separate physical-space phenomenon: quantized circulation}

Let a neutral condensate order parameter be written as \cite{Tsubota}
\begin{equation}
\psi(\bm{r})=\sqrt{\rho(\bm{r})}\,e^{\ii S(\bm{r})},
\end{equation}
where \(S\) is a dimensionless phase. Away from zeros of \(\psi\), the mechanical momentum per constituent is
\begin{equation}
\bm{p}=m\bm{v}=\hbar\nabla S.
\label{eq:phase-momentum}
\end{equation}
Single-valuedness of the order parameter implies
\begin{equation}
\Delta S=2\pi n,
\qquad n\in\mathbb{Z},
\end{equation}
and therefore
\begin{equation}
\boxed{
\oint_C \bm{p}\cdot\dd\bm{\ell}
=2\pi n\hbar
=nh.
}
\label{eq:circulation}
\end{equation}
Equivalently, the velocity circulation is \(nh/m\).

For an ideal straight vortex line, viewed in a two-dimensional transverse cross-section with core position \(\bm{r}_0\), Eq.~\eqref{eq:circulation} may be represented distributionally as
\begin{equation}
\boxed{
\nabla\times\bm{p}
=nh\,\delta^{(2)}(\bm{r}-\bm{r}_0)\,\hat{\bm{n}}.
}
\label{eq:vorticity}
\end{equation}
This formula is singular, topological, and system-dependent. It is not a uniform curl identity and does not concern the mixed derivatives \(\partial_x\) and \(\partial_p\) on phase space.

For a charged condensate, the gauge-invariant mechanical momentum takes the form
\begin{equation}
\bm{p}_{\mathrm{mech}}
=
\hbar\nabla S-q\bm{A},
\end{equation}
so the curl also contains the magnetic contribution \(-q\bm{B}\). This qualification prevents the neutral-superfluid formula from being applied outside its domain.
